% !TeX root = ../main.tex
\section{引言与问题分析}\label{sec:introduction}

\subsection{研究背景与灾害链特征}\label{sec:background}
% 简要介绍高寒山区冰岩崩—碎屑流—泥石流—山洪的链式灾害，以及吉隆事件的特殊性。
先介绍冰岩崩如何演变为碎屑流、泥石流和山洪，再结合文献说明现有研究能力，最后引出站点远离源区、资料尺度不同、精确计时困难等问题。文献按研究方向综合，不逐篇罗列。

\subsection{研究任务与关键难点}\label{sec:tasks}
% 不必照抄题目，而是将四问分别转化为数学问题，说明每问要解决的对象、关键难点和输出。
四问各写一段，说明任务、数据困难、解决思路和规定输出。第一问侧重关联排序，第二问研究触发边界，第三问研究时间区间与预警，第四问研究传播过程。

\subsection{总体技术路线}\label{sec:workflow}
% 给出贯穿四问的建模框架图，突出数据、模型和结果之间的传递关系。
% 框架图完成后保存在 figures/，再使用 figure 环境插入并添加 label。
用一张图展示数据、模型和输出的传递。体现空间资料整理与第一问可以并行，第四问的传播历时可以反馈第五章的到达预警。
